% P10-4a, 14.09.2026: Mechanismus, Quellenreichweite und v4-Nachweis angeglichen.
% M09, 11.09.2026: Leserführung und MCR konzentriert; neue
% HERKUNFT-M09-Kommentare unten sind die aktuelle Schreibgrundlage.
% Die folgende Versionshistorie beschreibt frühere Fassungen.
% M03 / B-18, 10.09.2026: Rückkopplung und Wiederverwendung
% am bestehenden Ablauf präzisiert; keine Änderung des Systems.
% M01 / B-09, 10.09.2026: Aufrufformen und Executor-Einbindung
% präzisiert; folgende Versionsnotizen dokumentieren den Vorzustand.
% ============================================================
% §6.2 Agent, Executor und Workflowgraph — v04 (Claude, 04.09.)
% v04 (04.09., v23-Lesungs-Runde): „Präge-" aus der Executor-
%   Aufzählung gestrichen (Kollision mit modellgestütztem „Claims
%   prägen", Kap. 4); Schlussrunde davor: „maßgeblicher
%   Projektzustand".
% v03/v04-Sammel (03./04.09.): V6-Zweistufigkeit; Gesamtpass
%   (Modellaufruf-Scope); Freeze-Pass (Instrumentierungs-
%   Anbindung); ANALYST-Spannweite mit Kollegen-Mikrofix „unter
%   anderem deterministisch" (disk-analyst.md §4-④). Danach
%   wieder SEMANTISCH FROZEN.
% v02 (Kollegen-Runde 1, alle 5 Punkte + Prüfauftrag gegengeprüft):
% ① „drei Komponentenarten" -> „zwei Knotenarten + Komposition"
% ② P0: „Kompositionsfehler scheitern laut" auf TYP-Verträge
%   begrenzt (kollidierte mit R-38!) + 6.5-Grenzen-Vorverweis
% ③ P0: Structured-Output entschärft (nur betrachtete Workflowstufen;
%   Schema VERLANGT Strukturen, erzwingt keine semantische
%   Stabilität; ID-Fallback code-verifiziert :215f)
% ④ GateLoop-Prosa ans Listing angepasst (Repair loopt, Rest ->
%   Finalisierung/Mensch; MaxAttempts/Historie ehrlich als „nicht
%   dargestellt")
% ⑤ CONTRADICT auf Proposal-before-Truth; „stufenweiser
%   Direktaufruf" -> „separat gestarteter Stufenworkflow"
% + Root-Input-Abgleich aufgelöst: Reproject == Steward-Recovery
%   (kein fünfter Typ). Kollegen-Urteil danach: semantisch
%   einfrieren, weiter mit 6.3.
% Grundlage: kapitel-6-gesamtvertrag.md ★FROZEN (§4 Zeile 6.2,
% Kernaussage #1 mit GateLoop integriert, 7+1-Ops-Satz §8.1) +
% probeschnitt-vertrag-6-1-6-2.md ★FROZEN (§3 Klassifikation v02,
% §4b Ingest-Pfad-Code-Walk, Nachtrag Ledger-Bestellung aus der
% 6.1-Runde).
% Form (§7): Rückverweis -> Instanziierung -> repräsentativer
% Mechanismus -> Beleg/Grenze. KEINE erneute Architekturlehre.
% Listing-Entscheid (Prüffrage bestanden): EIN Listing — die
% GateLoop-Kantenkomposition (zeigt die Systemgrenze „semantisches
% Urteil wird typisierte Nachricht, Routing ist deterministisch"
% besser als Prosa; L2/Enum verworfen — Kollegen-Einschätzung).
% Zählt 1 von max. 3–4 (Gesamtvertrag §7).
% TODO (Autor, Overleaf-Präambel): \usepackage{listings} + einmalige
% \lstset-Grundkonfiguration (bisher kein Listing in der Arbeit).
% Bewusst NICHT hier: Gate-/Responder-Mechanik (6.4) ·
% Kapsel-/Checkpoint-Details (6.5) · Speichernaht (6.7) ·
% Op-Katalog (nur 7+1-Satz + 2 Beispiele) · Placement-/Decision-/
% PBI-/Arch-Pfade (formgleich, EINMAL-Prinzip).
% ============================================================

\section{Agent, Executor und Workflowgraph}
\label{sec:realisierung-agent-executor-workflow}

Die Verantwortungsverteilung aus
Kapitel~\ref{chap:loesungsarchitektur} wird in MAF durch
\emph{Executors} und ihre Verbindung im Workflowgraphen umgesetzt.
Ein Executor bindet eine ausführbare Verarbeitungseinheit ein;
deren Aufgabe kann deterministisch oder modellgestützt sein. Ob
sie als \emph{agentisch} eingeordnet wird, richtet sich nach der
funktionalen Arbeitsdefinition in
Abschnitt~\ref{sec:arbeitsdefinition-llm-agent}, nicht allein nach
der verwendeten Frameworkklasse.

Die modellgestützten Stufen nutzen zwei technische Aufrufformen:
Der Ledger-Extraktor verarbeitet vorgegebenen Kontext über
\texttt{IChatClient.GetResponseAsync}; der Einordnungsagent nutzt
\texttt{AIAgent.RunAsync} und wählt Werkzeuge zur Bestandsabfrage
und Planablage. Beide Aufrufformen nutzen die gemeinsame
\emph{Chat-Pipeline}. Diese Verarbeitungskette verbindet den
Modellzugriff mit der Werkzeugausführung und einer konfigurierbaren
Protokollierung. Die Abschnitte~\ref{sec:realisierung-tools-mcp}
und~\ref{sec:realisierung-observability} erklären diese Anbindung.

Bei aktiviertem Structured Output fordert der Extraktor ein
JSON-Schema für Claims mit
Kennungen und Quellenankern an. Die Antwort wird in fachliche
Typen überführt; fehlende Kennungen werden deterministisch
ergänzt. Das Schema begrenzt die Ausgabeform, garantiert aber
keine semantisch richtige Zuordnung. Die vollständige
Verarbeitungsschrittfolge zeigt Tabelle~\ref{t:ledger-schritte}.
% HERKUNFT M09: SemanticLedgerExtractor.ExtractAsync und
% IngestionHitlResolveExecutor; M01-Abgrenzung IChatClient/AIAgent
% erhalten. Gemeinsame Instrumentierung: AgentChatPipelineBuilder.

Fachlich weiterverarbeitete Ergebnisse werden als typisierte
Nachrichten oder über Speicherwerkzeuge mit validierten
Argumenten übergeben. \emph{Deterministische Executors} übernehmen
unter anderem formale Prüfung, Verzweigung und Zustandsanwendung.
Die deklarierten Eingangs-, Sende- und Ausgabetypen erlauben
Kontrollen beim Graphaufbau und zur Laufzeit. Sie sichern die
technische Übergabe, nicht die Bedeutung des Ergebnisses oder
jede mögliche Graphverdrahtung
(Abschnitt~\ref{sec:realisierung-durabilitaet}).

\emph{Workflowgraphen} verbinden die Executors über Nachrichten
und Kantenbedingungen. Im Gesamtgraphen prüft der
\emph{Eingangs-Dispatcher}, welche von vier Eingaben vorliegt:
Transkript, strukturiertes Eingangspaket (Delta), Klärungsantworten
oder erneute Projektion des Bestands. Er gibt genau den
zugehörigen Nachrichtentyp weiter. Transkripte durchlaufen Ledger
und fachliche Ableitung; vorbereitete Deltas gelangen direkt zur
Verzweigung zwischen Initialisierung und Betrieb. Bei vorhandenem
Core nutzen beide Eingänge dieselben nachgelagerten Einordnungs-
und Fortschreibungsknoten. Wiederverwendet werden damit
Komponenten und ihre Verträge; eine dauerhafte Agentidentität
über mehrere Läufe ist dafür nicht erforderlich.

Ein frischer Lauf beginnt am Dispatcher. Die Wiederaufnahme eines
gesicherten Laufs am Checkpoint und die Einbettung von
Teilworkflows behandelt
Abschnitt~\ref{sec:realisierung-durabilitaet}.

\emph{Prüfung und Rückkopplung.} Die in
Abschnitt~\ref{sec:saeule-semantische-transformation} erklärte
Schleife lässt sich an der Einordnungsstufe konkret verfolgen.
Der Agent erhält Werkzeuge, um neue Inhalte, den Core und frühere
Ablehnungen abzufragen und einen Operationsplan zu speichern.
Ein deterministischer Checker prüft dessen formale Regeln.
Die gemeinsame Entscheidungsfunktion \texttt{GateLoop.Decide}
bestimmt aus Prüfergebnis, Reparierbarkeit und Versuchszahl den
Folgeschritt. Bei \texttt{Repair} erhält ein erneuter
Agentenaufruf die Fehlermeldungen; sein Ergebnis durchläuft
wieder den Checker. Versuchszähler und Historie führt die Stufe
in ihren typisierten Nachrichten mit.

Listing~\ref{lst:gateloop-komposition} zeigt diese Kanten im
separat startbaren Einordnungsworkflow. Alle anderen Ergebnisse
gelangen zur Finalisierung, die Plan und Prüfprotokoll ablegt.
Nur bei \texttt{Pass} und mindestens einer Operation
sendet sie eine Anfrage an den menschlichen Port. Die Statuswerte
\texttt{HumanReview} (manueller Prüfbedarf) und
\texttt{MaxAttemptsReached} (Versuchsschranke erreicht) führen hier
zum Abschluss ohne Anwendung. Der Prüfbedarf wird dokumentiert;
eine Portanfrage wird in diesen beiden Fällen nicht gestellt.
Ein leerer, bestandener Plan wird ohne menschliche Anfrage
durchgereicht. Dieser Leerfall sowie Versuchszähler und
Protokollierung sind im gekürzten Listing nicht dargestellt.
% HERKUNFT M09: IngestionResolveWorkflow.cs (Gate/Repair/History),
% IngestionHitlWorkflow.cs (Finalize: Pass + Ops > 0 -> Request;
% Pass + 0 Ops -> EmptyGate; übrige Entscheidungen -> YieldOutput).

% M12: Das kurze Listing mit seiner Beschriftung zusammenhalten.
\par\noindent
\begin{minipage}{\linewidth}
% M12-Nachtrag: Darstellung nur fuer dieses Listing, nicht in der Uni-Vorlage.
\lstset{
  language=[Sharp]C,
  basicstyle=\ttfamily,
  columns=fullflexible,
  keepspaces=true,
  showstringspaces=false,
  breaklines=true,
  breakatwhitespace=true
}
\begin{lstlisting}[caption={[Prüf-Reparatur-Kanten der Einordnungsstufe]Prüf-Reparatur-Kanten der Einordnungsstufe
  (gekürzter Ausschnitt aus \texttt{IngestionHitlWorkflow.cs}).},
  label={lst:gateloop-komposition}]
b.AddEdge(resolve, gate);
b.AddEdge<IngestionVerdict>(gate, repair,
    m => m is not null && m.Decision == GateDecision.Repair);
b.AddEdge<IngestionVerdict>(gate, finalize,
    m => m is not null && m.Decision != GateDecision.Repair);
b.AddEdge(repair, gate);
b.AddEdge(finalize, humanGate);
b.AddEdge(humanGate, apply);
\end{lstlisting}
\end{minipage}
\par\medskip

\emph{Modellurteil und formale Regel am Fall.} Der separate
Verfeinerungsfall F2 aus Abschnitt~\ref{sec:endtoend-pfade}
verdeutlicht die Aufgabenverteilung. Ein Diktat verlangte, dass
freigegebene Angehörige die Medikamenten-Einnahmen der letzten
sieben Tage einsehen können. Der Einordnungsagent bezog dies auf
die vorhandene, noch vage Anforderung REQ-42 zu Kalenderfunktion
und Medikamentengaben und schlug \texttt{REFINE} vor. Ob die
Aussage diese Anforderung fachlich verfeinert, ist das
Modellurteil. Die deterministische Zielregel prüft dagegen unter
anderem, ob ein bestehendes Zielelement angegeben ist; bei einer
Operation \texttt{NEW} wäre eine solche Zielangabe unzulässig.
Nach menschlicher Freigabe entstand eine neue Fassung von REQ-42
mit präzisiertem Zeitraum und Personenkreis; Vorgängerfassung und
Historie blieben erhalten. Der Plan bestand im ersten Prüfversuch,
eine Reparatur wurde in F2 nicht ausgelöst. Der Fall zeigt somit
Einordnung, Prüfung und autorisierte Anwendung, nicht die Wirkung
der Rückschleife. Die Gate-Interaktion und Speicherung vertiefen
die Abschnitte~\ref{sec:realisierung-human-gates} und
\ref{sec:realisierung-persistenz-aussenkopplung}.
% HERKUNFT M09: F2 20260817_123342_f23f06, plan.json (AF-1 ->
% REFINE REQ-42), ingest-gate-decisions.json, ingestion-attempts.json
% (ein Maker-Pass), applied/run-report.json und Core vor/nach Apply.

Die Einordnungs-Executors werden mit einer fachlichen
Konfiguration, dem \emph{Aspektprofil}, sowohl für Anforderungen
als auch für Architekturinformationen verwendet. Auch Einzel-
und Gesamtworkflow nutzen diese Verarbeitungskomponenten.
Ihre Kanten werden für diese Stufe jedoch an zwei Stellen
definiert: im Einzelworkflow und in der Komposition des
Gesamtgraphen. Die Wiederverwendung betrifft hier die
Ausführungslogik, nicht eine gemeinsame Kantenquelle.
% HERKUNFT M09: IngestionHitlWorkflow.Build versus
% PipelineComposedWorkflow.AddTo, Aufruf in PipelineFullWorkflow.
% Frühere Behauptung einer einzigen Quelle für DIESE Stufe korrigiert.

Auch der offenere Core-Analyst bleibt orchestriert: Ein eigener
Graph verteilt die Analyse auf mehrere Perspektiven, führt die
Funde deterministisch zusammen und gleicht sie gegen bekannte
Inhalte und frühere Ablehnungen ab. Ein modellgestützter Kritiker
bewertet die verbleibenden Kandidaten. Der Graph stellt die
beibehaltenen Funde als Delta für eine anschließende
Fortschreibung bereit; er schreibt sie nicht selbst in den Core.
Die verfügbaren Werkzeuge und ihre Anbindung behandelt der
folgende Abschnitt.
